
\subsection{Electroencefalografía (EEG)}\label{sec:eeg}

La electroencefalografía (EEG) consiste en la captación de la actividad eléctrica cerebral generada por la comunicación entre neuronas. Estas señales se registran mediante electrodos no invasivos colocados sobre el cuero cabelludo y son posteriormente amplificadas y representadas en función del tiempo \cite{luck_IntroductionEventRelated2014,sherman_EEGSignal2020}.

Dentro de estas señales se encuentran respuestas neuronales asociadas a distintos tipos de eventos (sensoriales, cognitivos y motores) conocidas como potenciales relacionados con eventos (\acrshort{erp}), así como aquellas asociadas a la imaginación motora (\acrshort{mi}) \cite{luck_IntroductionEventRelated2014,jiang_ComprehensiveReview2026}.

Es común realizar una separación de estas señales en función de su frecuencia, ya que se ha demostrado que distintas bandas de frecuencia se corresponden con estados mentales concretos \cite{sherman_EEGSignal2020,roohi-azizi_ChangesBrains2017}.

\begin{itemize}
    \item \textbf{Delta} (\textless4 Hz): Ubicadas en la región frontal en adultos y posterior en los niños. Se asocian principalmente al sueño, siendo también frecuentes en niños y en tareas de atención sostenida.

    \item \textbf{Theta} (4--7 Hz): Predominan en las regiones parietal y temporal en niños. En adultos, aparecen durante algunas fases del sueño o en estados de estrés

    \item \textbf{Alpha} (8--13 Hz): Se localizan en la región occipital. Se asocian a estados de relajación con los ojos cerrados, disminuyendo durante la concentración y el sueño.

    \item \textbf{Beta} (13--30 Hz): Son detectadas en los lóbulos frontal y parietal y presentan una distribución relativamente simétrica. Aparecen en tareas de calma activa con los ojos abiertos. Pero también pueden aparecer durante actividad mental intensa o situaciones de estrés.

    \item \textbf{Gamma} (\textgreater30 Hz): Se han relacionado con actividad en la región somatosensorial durante tareas de procesamiento multisensorial y en procesos cognitivos complejos, como la memoria.

    \item \textbf{Mu} (8--12 Hz): Se localizan en el córtex sensoriomotor y se asocian a la actividad de este sistema en reposo.
\end{itemize}

El análisis de señales \acrshort{eeg} es una herramienta muy potente que permite revelar información sobre la actividad cerebral de un individuo, teniendo muchas aplicaciones en neurociencia y en el control de \acrshort{bci} \cite{craik_DeepLearning2019}.

Entre las principales ventajas del uso de \acrshort{eeg} se encuentran \cite{sherman_EEGSignal2020,luck_IntroductionEventRelated2014, telpaz_UsingEEG2015,schirrmeister_DeepLearning2017}:

\begin{itemize}
    \item \textbf{Alta resolución temporal}: Permite registrar cambios en la actividad cerebral en el rango de milisegundos.

          \item\textbf{Medición directa de actividad eléctrica}: A diferencia de otras técnicas que estiman la actividad cerebral de forma indirecta mediante cambios en los niveles de oxígeno en sangre, las señales \acrshort{eeg} registran directamente la actividad eléctrica generada por poblaciones neuronales.

    \item \textbf{Carácter no invasivo}: La señal puede registrarse mediante electrodos colocados sobre el cuero cabelludo, sin necesidad de procedimientos quirúrgicos. Esto reduce el riesgo de sufrir complicaciones durante la captura de la señal y las molestias a los pacientes.

    \item \textbf{Bajo coste relativo}: En comparación con otras técnicas de neuroimagen, como la resonancia magnética funcional, los sistemas \acrshort{eeg} suelen ser más accesibles y fáciles de utilizar.

\end{itemize}

Algunas de las desventajas que convierten a estas señales en difíciles de analizar son las siguientes \cite{luck_IntroductionEventRelated2014,saha_IntraIntersubject2020,klonowski_EverythingYou2009}:


\begin{itemize}
    \item \textbf{Baja resolución espacial}: Las señales \acrshort{eeg} no permiten localizar con precisión la fuente de actividad cerebral, ya que representan una media de multiples fuentes de actividad neuronales.

    \item \textbf{Variabilidad inter-sujeto e intra-sujeto}: Existe una gran variabilidad en las señales EEG entre diferentes individuos, lo que dificulta la observación de patrones en ellas. A su vez, las propiedades estadísticas de las señales pueden variar a lo largo del tiempo (son señales no estacionarias), incluso para un mismo sujeto y tarea.

    \item \textbf{Ruido}: Las señales EEG son altamente susceptibles a artefactos, como el movimiento, la actividad muscular (EMG) o interferencias externas. Además, la amplitud de las señales EEG es muy baja en comparación con el ruido presente, lo que exige un preprocesado avanzado para eliminar el ruido y mejorar la calidad de las señales.

\end{itemize}

Estas limitaciones justifican el uso de técnicas avanzadas de procesamiento y aprendizaje automático. En los últimos años el uso de modelos de aprendizaje profundo (deep learning) ha supuesto un avance significativo, facilitando el desarrollo de aplicaciones más prácticas y reduciendo la dependencia de procesamiento manual por parte de expertos \cite{craik_DeepLearning2019}.

En el contexto del aprendizaje profundo (\textit{deep learning}) surge una cuarta limitación que es la \textbf{escasez de datos} disponibles para el entrenamiento. Los algoritmos de deep learning requieren una grán cantidad de muestras para entrenar y generalizar, lo que dificulta su aplicación en tareas con pocos ejemplos.

A diferencia de otros dominios donde pueden obtenerse grandes volúmenes de datos de forma relativamente sencilla, la adquisición de señales EEG requiere equipamiento especializado, personal cualificado y una preparación cuidadosa de cada sesión experimental. Además, la anotación de los datos suele ser costosa y dependiente del protocolo utilizado \cite{tan_DeepTransfer2018}.

La recogida de datos EEG también presenta limitaciones logísticas, ya que las sesiones experimentales suelen requerir tiempos prolongados de preparación y ejecución, lo que dificulta la obtención de medidas consecutivas \cite{sugden_RemoteCollection2023}. En aplicaciones BCI, este problema se ve agravado por la fatiga física y mental de los participantes, especialmente cuando se requieren múltiples ensayos o sesiones de larga duración \cite{rommel_DataAugmentation2022}.

Como consecuencia de estas limitaciones, los conjuntos de datos EEG suelen ser significativamente más reducidos que los utilizados habitualmente en otras áreas del aprendizaje profundo.
